% This document was generated by an LLM.
% Supplementary prerequisite review, not extracted from Hunter's lecture notes.

\documentclass{self-check}

\hypersetup{pdftitle={Self-Check Questions: Introductory Analysis Supplement}}

\title{Self-Check Questions:\\Introductory Analysis Supplement}
\subtitle{Prerequisite review: real sequences}
\author{}
\date{}

\begin{document}
\maketitle

\textit{Scope.} These supplementary questions review background assumed by
John K. Hunter's \textit{An Introduction to Real Analysis} (2012), especially
Chapter 1. They were not extracted from that lecture chapter and are separate
from the source-based questions in \texttt{intro-analysis-self-check.tex}.

\textit{How to use these questions.} Recall definitions and basic properties,
including their hypotheses and quantifiers. Full proofs are not required for
this first pass. No solutions are included.

All sequence terms and finite limits are real numbers unless specified
otherwise. Use \(\N=\{1,2,3,\ldots\}\).

\section{Sequences and convergence}
\begin{enumerate}
  \item Define a real sequence as a function. How does a sequence differ
    from the set of its values? May a sequence repeat a value?

  \item State the definition of \(x_n\to L\) using \(\epsilon\) and
    \(N\), with all quantifiers in order.
    What does it mean for a condition to hold eventually?

  \item In the definition of convergence, may \(N\) depend on
    \(\epsilon\)? May it depend on the index \(n\) quantified after it?
    Must one value of \(N\) work for every \(\epsilon>0\)?

  \item Write the precise negation of \(x_n\to L\) for a specified
    \(L\in\R\). How does this differ from saying that the sequence
    has no finite limit?

  \item Can a real sequence have two different finite limits?

  \item Does changing or deleting finitely many initial terms affect
    convergence or the value of a limit?

  \item Define a bounded sequence using an absolute-value inequality.
    Is it enough for all sufficiently late terms to satisfy such an
    inequality?

  \item Is every convergent real sequence bounded? Is every bounded
    real sequence convergent? Recall a simple counterexample to any
    implication that fails.
\end{enumerate}

\section{Subsequences}
\begin{enumerate}[resume]
  \item Define a subsequence \((x_{n_k})\) of \((x_n)\).
    What conditions must the indices \(n_k\) satisfy?
    Must \(n_k\to\infty\)? May equal values occur in a subsequence?

  \item If \(x_n\to L\), what can be said about every subsequence?
    Must all terms of a subsequence be different from \(L\)?

  \item Does the existence of one convergent subsequence imply that
    the original sequence converges? Recall a simple counterexample
    if it does not.

  \item What can be concluded if a sequence has two subsequences
    converging to different finite limits?
\end{enumerate}

\section{Basic limit properties}
\begin{enumerate}[resume]
  \item Suppose \(x_n\to x\), \(y_n\to y\), and \(c\in\R\).
    State the limit rules for \(x_n+y_n\), \(cx_n\), \(x_ny_n\),
    and \(x_n/y_n\). What extra hypotheses are needed for the quotient?

  \item Suppose \(x_n\to x\), \(y_n\to y\), and eventually
    \(x_n\le y_n\). What inequality follows for the limits?
    If eventually \(x_n<y_n\), must the limit inequality be strict?

  \item State the squeeze (or sandwich) theorem for real sequences,
    including the required inequalities and limits.

  \item Define nondecreasing and nonincreasing sequences.
    How do these differ from strictly increasing and strictly decreasing
    sequences?

  \item State the monotone convergence theorem for real sequences.
    Which one-sided bound is needed in each case?
    Express the limit as a supremum or infimum of the sequence's values.

  \item Define \(x_n\to+\infty\) and \(x_n\to-\infty\) using
    quantifiers. Does every sequence with no finite limit diverge to
    one of these infinities? Recall a simple counterexample if it does not.
\end{enumerate}

\end{document}
